\documentclass{article}
% NeurIPS 2022 format, final (non-anonymous) mode, as for the proposal.
\PassOptionsToPackage{numbers}{natbib}
\usepackage[preprint]{neurips_2022}
\usepackage[utf8]{inputenc}
\usepackage[T1]{fontenc}
\usepackage{hyperref}
\usepackage{url}
\usepackage{booktabs}
\usepackage{amsfonts}
\usepackage{amsmath}
\usepackage{amssymb}
\usepackage{microtype}
\usepackage{xcolor}
\usepackage{enumitem}
\usepackage{listings}
\lstset{basicstyle=\ttfamily\footnotesize, columns=fullflexible, keepspaces=true, aboveskip=2pt, belowskip=2pt, xleftmargin=1em}

\hypersetup{colorlinks=true, linkcolor=blue!60!black, citecolor=blue!60!black, urlcolor=blue!60!black}
\setlength{\abovedisplayskip}{3pt}
\setlength{\belowdisplayskip}{3pt}
\setlist[itemize]{leftmargin=1.2em, itemsep=1pt, topsep=1pt, parsep=0pt}
\makeatletter
\renewcommand{\@toptitlebar}{\hrule height 4\p@ \vskip 0.18in \vskip -\parskip}
\renewcommand{\@bottomtitlebar}{\vskip 0.18in \vskip -\parskip \hrule height 1\p@ \vskip 0.02in}
\makeatother
\newcommand{\code}[1]{\texttt{#1}}
\newcommand{\E}{\mathbb{E}}

\title{Learning the Graph for Relational Deep Learning\\[2pt]\large Milestone report}
\author{%
  Harshal Kulkarni \quad \texttt{hsk8171@stanford.edu} \\
  Project 3: PR to \texttt{torch\_geometric.contrib} + blog post
}

\begin{document}
\maketitle
\vspace{-2.6em}

\section{Problem and paper summary}
\vspace{-0.4em}
Relational deep learning~\cite{fey2024rdl} turns a database into a graph mechanically: rows
become nodes and foreign keys become edges. \emph{What Makes a Desired Graph for Relational Deep
Learning?}~\cite{cheng2026desired} argues that this schema graph is often poor for message
passing: irrelevant columns and tables add noise (\emph{information overload}), and related
entities sit many hops apart (\emph{semantic fragmentation}). Its structural optimizer learns the
graph together with the model, in four parts: \textbf{A}~filtering (learned switches on columns
and tables, plus a variational bottleneck), \textbf{B}~injection (new edges from KNN, two-hop,
behavioural and temporal templates), \textbf{C}~gated aggregation (a learned switch per added
relation), and \textbf{D}~an $L_0$ sparsity penalty. On RelBench~\cite{robinson2024relbench} it
reports, e.g., \code{driver-top3} ROC-AUC 81.5 (Base) $\to$ 84.7 (filtering only) $\to$ 86.8 (full).
We implement it as PyG \code{contrib} modules and test these claims. This milestone covers
\textbf{A} (switches and penalty) end to end: code, tests and a full experiment.

\section{Method: filtering (paper Eq.~3--5, 11)}
\vspace{-0.4em}
\textbf{Base model} (RelBench's relational GraphSAGE). Each table's encoder turns row $v$ into
$h^{(0)}_v$; two GraphSAGE layers and a linear head give the prediction $\hat y_v$.

\textbf{Gates.} Every column $c$ of table $T$ gets a gate $g_{T,c}$ and every table a gate $s_T$,
each with one learned logit $\theta$. In training a gate is a Hard-Concrete
sample~\cite{louizos2018l0}; at test time it is a hard 0/1 switch (Eq.~4):
\begin{align*}
\text{train: } g &= \min\!\Big(1, \max\!\big(0,\ \sigma\big(\tfrac{\log u - \log(1-u) + \theta}{t}\big)(\zeta-\gamma)+\gamma\big)\Big),\quad u \sim U(0,1);\\
\text{test: } g &= \mathbb{I}[\sigma(\theta) > \tau].
\end{align*}
The gates multiply what the model sees, and the loss pays for every gate left open (Eq.~11):
\begin{gather*}
\tilde e_{v,c} = g_{T,c}\, e_{v,c}, \qquad \tilde h^{(0)}_v = s_T\, h^{(0)}_v,\\
\mathcal L = \mathcal L_\text{task} + \lambda \sum_{T,c} \E\|g_{T,c}\|_0 + \lambda_s \sum_T \E\|s_T\|_0,
\qquad \E\|g\|_0 = \sigma\!\big(\theta - t\log\tfrac{-\gamma}{\zeta}\big).
\end{gather*}
Here $e_{v,c}$ is column $c$'s 128-d encoding; the column gate acts inside the encoder, before
columns are mixed, so a closed column is removed exactly. The paper's Eq.~3 gates the 128
dimensions of $h^{(0)}$ instead; we gate columns so the result is readable (``the model dropped
\code{results.rank}'') and document the difference. No table gate is put on the target table.
Settings: $\gamma=-0.1$, $\zeta=1.1$, $t=2/3$, $\tau=0.5$.

\section{What is built}
\vspace{-0.4em}
\begin{itemize}
\item \textbf{PyG fork} (\code{contrib}, merged after review): \code{HardConcreteGate} and
  \code{RelationalFilter}, with PyG-style API docstrings (equations, arguments, shapes, examples)
  and 13 unit tests. Usage, independent of RelBench:
\begin{lstlisting}[language=Python]
filt = RelationalFilter(num_columns={'results': 11, ...}, table_types=[...])
e = filt.gate_columns('results', e)      # [N, C, 128]: column c scaled by g_c
h_dict = filt.gate_tables(h_dict)        # each table's [N, 128] scaled by s_T
col_pen, tab_pen = filt.l0_penalties()   # loss += lambda*col_pen + lambda_s*tab_pen
\end{lstlisting}
\item \textbf{Experiment code}: one training script that reproduces RelBench's official example within
  seed noise; \code{--model filter}; a resumable sweep runner; a Colab notebook for GPU runs.
\item \textbf{Correctness checks.} With every gate fixed at 1 the gated model trains \emph{exactly}
  like Base (same loss at every step, same weights), and over 5 seeds it scores the same
  ($p=0.97$). A closed gate removes exactly its column or table: the prediction no longer depends
  on its values (35 experiment tests in all).
\item \textbf{Found along the way.} A bug in PyTorch Frame 0.3.0: a missing value gives its
  column's weights a NaN gradient, so mostly-missing columns stop learning after one step (fixed
  behind a flag). Encoding each sampled row once per batch makes training $2.5\times$ faster with
  the same accuracy.
\end{itemize}

\section{Experiment and results}
\vspace{-0.4em}
\textbf{Setup.} \code{rel-f1}, task \code{driver-top3} (will a driver finish top 3; 1{,}353 / 588 /
726 train / val / test rows), ROC-AUC. To avoid picking noise, the protocol was fixed in advance:
\textbf{(A)} tune Base on the paper's grid (learning rate $\{10^{-4},10^{-3},10^{-2}\}$ $\times$
hidden $\{64,128,256\}$, 3 seeds); \textbf{(B)} tune the filter with the same budget,
$\lambda,\lambda_s \in \{10^{-4},10^{-3},10^{-2}\}$, 3 seeds; \textbf{(C)} compare the two winners on
5 fresh seeds and read test once. All runs: 30 epochs, NaN fix on. Gate start $\theta_0=5$ and a
separate gate learning rate of 0.1 were chosen in 7 dry runs: with the model's learning rate a gate
can move only $\sim$0.3 in the 60--90 training steps this small task allows, and starting at
$\theta_0=2$ made gates noisy enough to cost $\sim$0.03 even when none closed.

\begin{center}
\vspace{-0.6em}
\small
\begin{tabular}{@{}lccc@{}}
\toprule
 & Validation (ours) & Test (ours) & Test (paper, Table 1) \\
\midrule
Base (lr $10^{-2}$, hidden 64) & $0.785 \pm 0.008$ & $0.766 \pm 0.040$ & $0.815 \pm 0.011$ \\
Base + filtering ($\lambda=\lambda_s=10^{-4}$) & $0.790 \pm 0.014$ & $0.772 \pm 0.041$ & $0.847 \pm 0.006$ \\
\midrule
Difference & $+0.005$ ($p=0.50$) & $+0.006$ ($p=0.81$) & $+0.032$ \\
\bottomrule
\end{tabular}
\vspace{-0.4em}
\end{center}

\textbf{Findings.}
\begin{itemize}
\item \textbf{No measurable gain from filtering} on this task: the paper's +3.2 points is not
  reproduced. Our test scores vary across seeds (0.72--0.81) far more than the paper reports.
\item \textbf{$\lambda$ controls sparsity}, as the paper's Theorem~3.2 needs: by epoch 30, 31, 25 and
  13 of 45 columns stay open for $\lambda = 10^{-4}, 10^{-3}, 10^{-2}$. With only 13 open,
  validation is still within one std of Base: most columns are redundant, but removing them does
  not help.
\item \textbf{The gates act too late.} They start closing around epoch 14, while the best
  validation epoch is around 16, so the selected models are barely filtered (mostly the 3 tables a
  2-layer model never reaches). Too fast a gate learning rate closes useful columns instead.
\item Our Base is about 5 points below the paper's Base even with the paper's tuning grid.
\end{itemize}

\section{Next steps}
\vspace{-0.4em}
\textbf{(1)} Learn the gates first, then freeze them at 0/1 before choosing the final model.
\textbf{(2)} Injection (KNN, two-hop, behavioural, temporal edges) with gated aggregation (B, C) and
the bottleneck (D), then the full optimizer against the paper's 0.868. \textbf{(3)} More tasks
(\code{driver-dnf}, \code{driver-position}, \code{rel-trial}). \textbf{(4)} PR polish and the blog post.
Code: \href{https://github.com/HARSHALK2598/pytorch_geometric/tree/contrib/structural-optimizer}{PyG fork,
branch \code{contrib/structural-optimizer}}; experiments, results and a run-by-run lab notebook:
\href{https://github.com/HARSHALK2598/cs224w-structural-optimizer}{project repository}.

\footnotesize
\begin{thebibliography}{9}
\setlength{\itemsep}{0pt}
\bibitem{cheng2026desired} Y.~Cheng and S.~Luo. What makes a desired graph for relational deep learning? In \emph{ICML}, 2026. arXiv:2606.08491.
\bibitem{fey2024rdl} M.~Fey, W.~Hu, K.~Huang, J.~E. Lenssen, R.~Ranjan, J.~Robinson, R.~Ying, J.~You, and J.~Leskovec. Position: Relational deep learning: Graph representation learning on relational databases. In \emph{ICML}, 2024.
\bibitem{robinson2024relbench} J.~Robinson, R.~Ranjan, W.~Hu, K.~Huang, J.~Han, A.~Dobles, M.~Fey, J.~E. Lenssen, Y.~Yuan, Z.~Zhang, X.~He, and J.~Leskovec. RelBench: A benchmark for deep learning on relational databases. In \emph{NeurIPS Datasets and Benchmarks}, 2024.
\bibitem{louizos2018l0} C.~Louizos, M.~Welling, and D.~P. Kingma. Learning sparse neural networks through $L_0$ regularization. In \emph{ICLR}, 2018.
\end{thebibliography}

\end{document}
